%
% spec-hook-gates.tex
%
% Composed Hook Gate Machine: A Z Specification
% A formal model of two hook decisions implemented in
% internal/hook (subagent_start.go, pretooluse.go), current as of
% the commit on branch feat/z-spec-state-machines that added this
% document, in support of bead ethos-mr1z (workstream 3 of the
% formal-spec gap program), formalizing the observability gap
% ethos-yf6n identifies.
%
% Compile: pdflatex spec-hook-gates.tex
%
% Must be .tex, not .md -- see spec-mission-lifecycle.tex's own
% header comment for why: fuzz and pdflatex do not care about the
% extension, but this repo's z-spec tooling globs for .tex.
%

\documentclass[a4paper,10pt,fleqn]{article}
\usepackage[margin=1in]{geometry}
\usepackage{fuzz}
\usepackage{booktabs}
\usepackage{tabularx}
\usepackage[colorlinks=true,linkcolor=blue,citecolor=blue,urlcolor=blue]{hyperref}

\newcommand{\fpath}[1]{\texttt{#1}}
\newcommand{\field}[1]{\texttt{#1}}

\begin{document}

\title{Composed Hook Gate Machine: A Z Specification}
\author{Mike S, for Punt Labs}
\date{September 2026}
\hypersetup{
  pdftitle={Composed Hook Gate Machine: A Z Specification},
  pdfauthor={Punt Labs},
  pdfsubject={Z Specification},
  pdfcreator={fuzz/probcli}
}
\maketitle
\tableofcontents
\newpage

% ===================================================================
\section{Scope}
\label{sec:scope}
% ===================================================================

This specification models two hook decisions as one composed
machine, because the second is a function of the first and no
document that treated them separately could state the fact that
matters: a single failure at the first decision silently disables
both.

\begin{enumerate}
  \item The \textbf{SubagentStart verifier gate}
    (\field{checkVerifierHash}, subagent\_start.go:585--652), which
    decides whether a freshly spawned subagent is treated as the
    frozen evaluator of an open mission, and, if so, replaces its
    persona context with the isolation block and sets
    \field{ETHOS\_VERIFIER\_ALLOWLIST} in its environment
    (\field{buildVerifierIsolationBlock}/\field{buildVerifierAllowlistEnv},
    subagent\_start.go:178--194).
  \item The \textbf{PreToolUse write gate}
    (\fpath{internal/hook/pretooluse.go}:110--128), which decides
    whether a \field{Write} or \field{Edit} tool call inside that
    subagent's session is confined to the mission's write set. Its
    only input is the environment variable the first gate sets.
\end{enumerate}

The bug class this document exists to make checkable is not a wrong
answer at either gate individually --- both gates are separately
correct on the inputs they receive. It is a \emph{silent} loss of
input at the seam between them: \field{MISSION\_ID}, the environment
variable that tells the SubagentStart gate which mission (if any)
this spawn serves, is set by a Tier B \fpath{PreToolUse} dispatch
(pretooluse\_dispatch.go:894) that runs in a \emph{different} process
invocation, at spawn time, before the subagent the gate is deciding
about even exists. If that variable does not arrive --- for any
reason --- \field{checkVerifierHash} returns \field{nil, nil}
today, at subagent\_start.go:599--605, exactly as if this spawn had
nothing to do with any mission at all. The two invariants above are
not violated; they are never evaluated. Bead ethos-yf6n is triaged
\textbf{DO NOT FIX AS PROPOSED} (2026-07-31): the code behaves as
designed, and the gap is observability, not gating semantics. This
document exists to say precisely what ``the gap is observability''
means, calculate the smallest change that closes it, and record why
that smaller change is sufficient.

\subsection{Modelling choice: one spawn at a time}
\label{sec:onespawn}

\field{HandleSubagentStartWithDeps} decides once per subagent
spawn (subagent\_start.go:102--235); \field{HandlePreToolUse}
decides once per tool call inside that spawn's session
(pretooluse.go:86--163). This specification's state,
\field{GateState} (\S\ref{sec:state}), describes one spawn's
worth of that history: the SubagentStart decision, taken exactly
once, followed by zero or more PreToolUse write decisions that all
read the same environment the first decision set. Nothing here would
change if it were promoted to a partial function from session
identifiers to \field{GateState} --- exactly the same simplification
\S1.1 of \fpath{spec-mission-lifecycle.tex} makes for missions, for
the same reason: every decision in scope touches exactly one spawn's
state.

\subsection{Non-goals}
\label{sec:nongoals}

The following are real mechanisms in the two files this document
models, deliberately excluded because none of them bears on the
silent-disabled property:

\begin{itemize}
  \item \textbf{\field{agentType} empty} (subagent\_start.go:596--598)
    and \textbf{no mission store configured}
    (\field{deps.Missions == nil}, subagent\_start.go:586--588).
    Neither can produce a false ``not a verifier'' --- with no agent
    type there is no evaluator handle to match against, and with no
    store there is no notion of an open mission at all. Both return
    \field{nil, nil}, but from a spawn that provably has no verifier
    role to lose.
  \item \textbf{Hash-source misconfiguration}
    (\field{deps.Hash.Validate()} failing, subagent\_start.go:589--595)
    and \textbf{evaluator hash mismatch}
    (subagent\_start.go:645--650). Both \emph{refuse} the spawn
    outright (a non-nil error, a fatal launch failure the operator
    sees) rather than silently falling through ungated. They are the
    opposite failure mode from the one in scope and are excluded for
    the same reason \fpath{spec-mission-lifecycle.tex} excludes
    \field{Abandon}'s \field{reason} argument: a real mechanism with
    no bearing on the invariant at hand.
  \item \textbf{\field{ETHOS\_VERIFIER\_EXTRACT\_INTO}} (DES-052,
    pretooluse.go:130--150) and \textbf{path-prefix matching}
    (\field{splitAllowlist}/\field{pathAllowed}). The mission brief
    for this document asks that \field{ETHOS\_VERIFIER\_ALLOWLIST} be
    modelled as \emph{the} input to the write gate; the extract-into
    carve-out is a second, independent env var this document does
    not model, and the prefix-matching mechanics are abstracted to a
    single boolean, \field{targetInAllowlist?}, in
    \S\ref{sec:writeops}.
  \item \textbf{\field{denyInRepoMCPWrite}} (pretooluse.go:115--117)
    and the \textbf{Tier B preconditions admission gate}
    (\field{evalContractPreconditions}, pretooluse.go:106--108). Both
    are additional deny paths that fire before the allowlist check,
    on different grounds (an in-repo MCP write tool class; an
    unsatisfied must-read precondition). Neither one's absence is
    what the silent-disabled property is about.
  \item \textbf{Tier A advice} (\field{dispatchAgent}'s
    non-\field{MISSION\_ID} branches, pretooluse\_dispatch.go:80--132)
    is orthogonal by the code's own comment (pretooluse.go:69):
    advice fires on \field{Agent} tool calls, which never carry a
    \field{file\_path}; the write gate fires on \field{Write}/
    \field{Edit}. The two hooks share one payload read, not one
    decision.
\end{itemize}

% ===================================================================
\section{Basic Types and Free Types}
\label{sec:basic}
% ===================================================================

One given set, carrying no internal structure relevant to this
model: mission identifiers.

\begin{zed}
[MISSIONID]
\end{zed}

fuzz has no built-in boolean, and \field{BOOL}/\field{true}/
\field{false} collide with B keywords under ProB's translation, so
this document uses the same convention as
\fpath{spec-mission-lifecycle.tex}:

\begin{zed}
ZBOOL ::= ztrue | zfalse
\end{zed}

\subsection{The MISSION\_ID channel}
\label{sec:channel}

The value \field{os.Getenv("MISSION\_ID")} reads at
subagent\_start.go:128 is a Go string, and Go strings have no
``absent'' distinct from \field{""}. This specification gives the
channel the explicit third value Go's type does not: a declared
mission identifier, or an explicit absence, are different values of
one free type, not a present string that happens to be empty.

\begin{zed}
MissionIDChannel ::= midAbsent | midPresent \ldata MISSIONID \rdata
\end{zed}

\field{channel? = midAbsent} is exactly the guard
\field{strings.TrimSpace(declaredMissionID) == ""} at
subagent\_start.go:600.

\subsection{Why the channel goes empty}
\label{sec:causes}

Three distinct real-world events all reach the gate as
\field{channel? = midAbsent}, and the mission brief for this
document asks which of them the gate can tell apart:

\begin{zed}
AbsenceCause ::= causeTierAAdHoc | causePreToolUseMisdispatch
               | causePreToolUseNotInstalled
\end{zed}

\begin{itemize}
  \item \field{causeTierAAdHoc} --- the spawn genuinely has no
    mission (\field{dispatchTierBOrTierA} fell through to Tier A,
    pretooluse\_dispatch.go:132, because no pending dispatch and no
    explicit \field{MISSION\_ID} matched). Absence here is correct:
    there is no mission for this spawn to be the verifier of.
  \item \field{causePreToolUseMisdispatch} --- a Tier B dispatch
    \emph{did} run and \emph{did} write
    \field{additionalEnv["MISSION\_ID"]} (pretooluse\_dispatch.go:894),
    but Claude Code's env propagation into the spawned subagent's
    process failed to carry it through by the time
    \field{HandleSubagentStartWithDeps} reads
    \field{os.Getenv("MISSION\_ID")} at subagent\_start.go:128.
  \item \field{causePreToolUseNotInstalled} --- the operator's Claude
    Code settings never wire \field{HandlePreToolUse} to the
    \field{PreToolUse} hook event at all, so \field{dispatchAgent}
    (pretooluse\_dispatch.go:80) never runs for this spawn, and no
    \field{MISSION\_ID} is ever written to inject.
\end{itemize}

All three reach subagent\_start.go:599--605 as the identical value
\field{channel? = midAbsent}. \S\ref{sec:theorems}'s
\field{CauseIndistinguishable} states, and the operation schemas in
\S\ref{sec:decision} demonstrate by construction, that none of the
three is distinguishable at this gate --- not because the current
code fails to check something it could, but because the channel
carries no signal that would let it.

% ===================================================================
\section{State}
\label{sec:state}
% ===================================================================

\begin{schema}{GateState}
verifierOfOpenMission : ZBOOL \\
decided : ZBOOL \\
gateApplied : ZBOOL \\
diagnosticEmitted : ZBOOL \\
allowlistActive : ZBOOL
\where
allowlistActive = gateApplied
\end{schema}

Five fields, one invariant:

\begin{itemize}
  \item \field{verifierOfOpenMission} is \textbf{ground truth}: does
    an open mission exist, right now, that names this spawn's
    \field{agentType} as its evaluator? This is a fact about the
    on-disk mission store and this spawn's real agent type. It does
    not depend on whether \field{MISSION\_ID} reached this hook
    invocation --- the mission is open, or it is not, whether or not
    anyone told this process its identifier. \S\ref{sec:init} fixes
    it once, at spawn time, and every operation in
    \S\ref{sec:decision}--\ref{sec:writeops} carries it through
    unchanged; \field{GroundTruthFrozen} (\S\ref{sec:theorems})
    records this.
  \item \field{decided} records whether the one-shot SubagentStart
    decision (\S\ref{sec:decision}) has fired yet for this spawn.
  \item \field{gateApplied} is the \emph{observable} counterpart to
    \field{verifierOfOpenMission}: did Phase 3.5's isolation block
    fire (\field{len(verifierMissions) > 0},
    subagent\_start.go:179)? When the channel carries a mission ID,
    this tracks ground truth exactly (\S\ref{sec:decision}'s guards
    enforce the agreement). When it does not, the gate has nothing to
    check against and \field{gateApplied} is unconditionally
    \field{zfalse} regardless of what \field{verifierOfOpenMission}
    says --- this is the whole bug, stated as a single field
    disagreeing with another.
  \item \field{diagnosticEmitted} records whether the decision wrote
    anything an operator could see. Under \emph{today's} code this is
    \field{zfalse} on every \field{midAbsent} decision
    (subagent\_start.go:599--605 returns bare \field{nil, nil}, no
    \field{fmt.Fprintf}); \S\ref{sec:validation}'s corrected variant
    changes exactly this one field's value on exactly that path.
  \item \field{allowlistActive} models
    \field{ETHOS\_VERIFIER\_ALLOWLIST} being set in the subagent's
    environment. The invariant
    \field{allowlistActive = gateApplied} is the \textbf{composition
    fact} this whole document exists to state formally:
    \field{buildVerifierAllowlistEnv} is called from inside the
    \field{len(verifierMissions) > 0} branch and nowhere else
    (subagent\_start.go:179--194), so the write gate's one input is
    set if and only if the SubagentStart gate applied. Nothing in
    \fpath{pretooluse.go} sets or reads
    \field{ETHOS\_VERIFIER\_ALLOWLIST} independently of that branch.
\end{itemize}

% ===================================================================
\section{Initialization}
\label{sec:init}
% ===================================================================

\field{Init} corresponds to the state of a freshly spawned subagent,
the instant before \field{HandleSubagentStartWithDeps} makes its
one decision: nothing has fired, nothing has been observed.

\begin{schema}{Init}
GateState'
\where
decided' = zfalse \\
gateApplied' = zfalse \\
diagnosticEmitted' = zfalse
\end{schema}

\field{verifierOfOpenMission}$'$ is left unconstrained beyond
\field{GateState}'s own type: whether this spawn really is the
evaluator of a real open mission is a fact about the world fixed
before the hook runs, not a choice \field{Init} makes. ProB
enumerates both values, exactly as \fpath{spec-mission-lifecycle.tex}
leaves \field{evaluatorHash}$'$ and \field{repoRootSet}$'$
unconstrained at \field{Init} for the same reason.
\field{allowlistActive}$'$ follows from the schema invariant
(\field{gateApplied}$'$ = \field{zfalse}, so
\field{allowlistActive}$'$ = \field{zfalse}): no tool call has
happened yet, so nothing has set the environment.

% ===================================================================
\section{Operations}
\label{sec:ops}
% ===================================================================

\subsection{The SubagentStart decision}
\label{sec:decision}

Exactly one of the four operations below fires per spawn, guarded by
\field{decided = zfalse}; each sets \field{decided}$'$ =
\field{ztrue} so none of the four --- nor any other decision
operation --- is enabled again for this spawn.
\field{SubagentStartOneShot} (\S\ref{sec:theorems}) records this.

\textbf{MuteSkip} models \emph{today's} code, subagent\_start.go:
599--605: the channel is empty, and \field{checkVerifierHash}
returns \field{nil, nil} without inspecting
\field{verifierOfOpenMission} at all and without writing anything an
operator could see. \field{cause?} is declared --- the mission brief
asks this document to carry it --- but appears nowhere in the
predicate, which is the point: this operation cannot look at its own
input to decide anything, because the real code has no such input
either.

\begin{schema}{MuteSkip}
\Delta GateState \\
channel? : MissionIDChannel \\
cause? : AbsenceCause
\where
decided = zfalse \\
channel? = midAbsent \\
decided' = ztrue \\
gateApplied' = zfalse \\
diagnosticEmitted' = zfalse \\
verifierOfOpenMission' = verifierOfOpenMission
\end{schema}

\textbf{NotAVerifier} models the two branches where the channel
\emph{does} carry a mission ID and \field{checkVerifierHash} correctly
determines this spawn is not that mission's verifier --- the mission
is not open (subagent\_start.go:611--615) or its evaluator is someone
else (subagent\_start.go:616--621, the ethos-z69l fix). Both return
\field{nil, nil} with no diagnostic, and both are safe: the gate
\emph{did} look, and ground truth agrees there was nothing to gate.
That agreement is written as an explicit guard, not derived, because
it is the fact that makes this branch harmless where
\field{MuteSkip} is not.

\begin{schema}{NotAVerifier}
\Delta GateState \\
channel? : MissionIDChannel \\
missionOpen?, isEvaluator? : ZBOOL
\where
decided = zfalse \\
channel? \neq midAbsent \\
\lnot (missionOpen? = ztrue \land isEvaluator? = ztrue) \\
verifierOfOpenMission = zfalse \\
decided' = ztrue \\
gateApplied' = zfalse \\
diagnosticEmitted' = zfalse \\
verifierOfOpenMission' = verifierOfOpenMission
\end{schema}

\textbf{VerifierGateHit} models the success path,
subagent\_start.go:611--651 falling through both checks: the mission
is open and this spawn is its evaluator. \field{diagWarn?} carries
the one remaining real distinction on this path --- a pre-3.3
mission with an empty pinned hash warns and skips the hash compare
(subagent\_start.go:624--635, \field{diagnosticEmitted}$'$ =
\field{ztrue}); a post-3.3 mission with a matching hash succeeds
silently (the comment at subagent\_start.go:101: ``Hash success is
silent''). A hash \emph{mismatch} refuses the spawn outright
(subagent\_start.go:645--650) and is out of scope by
\S\ref{sec:nongoals} --- it never reaches this operation.

\begin{schema}{VerifierGateHit}
\Delta GateState \\
channel? : MissionIDChannel \\
missionOpen?, isEvaluator?, diagWarn? : ZBOOL
\where
decided = zfalse \\
channel? \neq midAbsent \\
missionOpen? = ztrue \\
isEvaluator? = ztrue \\
verifierOfOpenMission = ztrue \\
decided' = ztrue \\
gateApplied' = ztrue \\
diagnosticEmitted' = diagWarn? \\
verifierOfOpenMission' = verifierOfOpenMission
\end{schema}

\subsection{The PreToolUse write decision}
\label{sec:writeops}

\field{ETHOS\_VERIFIER\_ALLOWLIST} is the one input this document
models for the write gate (\S\ref{sec:nongoals}); path-prefix
matching against it is abstracted to a single boolean,
\field{targetInAllowlist?}. All three operations below are guarded
by \field{decided = ztrue}: no tool call is dispatched to a subagent
before the SubagentStart decision that spawned it has run.

\textbf{WriteAllowedNoAllowlist} models pretooluse.go:110--113's
unconditional early return. It takes no \field{targetInAllowlist?}
input at all --- not merely an unconstrained one --- because the
real code never evaluates the target path on this branch; the
absence of the parameter is the formal witness that the decision is
vacuous over the tool call's target, the same technique
\S\ref{sec:decision} uses for \field{cause?}.

\begin{schema}{WriteAllowedNoAllowlist}
\Xi GateState
\where
decided = ztrue \\
allowlistActive = zfalse
\end{schema}

\begin{schema}{WriteAllowedInAllowlist}
\Xi GateState \\
targetInAllowlist? : ZBOOL
\where
decided = ztrue \\
allowlistActive = ztrue \\
targetInAllowlist? = ztrue
\end{schema}

\begin{schema}{WriteDenied}
\Xi GateState \\
targetInAllowlist? : ZBOOL
\where
decided = ztrue \\
allowlistActive = ztrue \\
targetInAllowlist? = zfalse
\end{schema}

None of the three changes any component of \field{GateState}: the
write gate consumes the SubagentStart decision, it does not revise
it. This is why \S\ref{sec:validation}'s reachability demonstration
needs only \S\ref{sec:decision}'s operations --- the write gate is a
theorem about the decision, not a further place the decision can go
wrong.

% ===================================================================
\section{Theorems}
\label{sec:theorems}
% ===================================================================

fuzz type-checks schemas; it does not discharge proof obligations.
The first two claims below are checked by ProB animation
(\S\ref{sec:validation}); the last three are backed by inspection of
the schemas above, in the same sense
\fpath{spec-mission-lifecycle.tex}'s \field{EvaluatorHashFrozen} is:
no reachable violation is even expressible, so animation would
confirm nothing beyond what reading the predicates already shows.

\begin{center}
\begin{tabularx}{\textwidth}{@{}lX@{}}
\toprule
\textbf{Theorem} & \textbf{Justification} \\
\midrule
\field{SilentDisabledReachableUnderMuteSkip} &
  Checked by ProB, not by inspection: \S\ref{sec:validation}'s
  decoupled model with \field{MuteSkip} in place of
  \field{LoudSkip} reaches, by \field{INITIALISATION} then one
  \field{MuteSkip} firing with \field{verifierOfOpenMission?} chosen
  \field{ztrue}, a state satisfying \field{decided = ztrue} $\land$
  \field{verifierOfOpenMission = ztrue} $\land$ \field{gateApplied =
  zfalse} $\land$ \field{diagnosticEmitted = zfalse} --- the silent-
  disabled state --- and reports it as a deadlock. \\
\field{SilentDisabledUnreachableUnderLoudSkip} &
  Checked by ProB: the same decoupled model with \field{LoudSkip} in
  place of \field{MuteSkip} explores every reachable state
  (\S\ref{sec:validation}) and finds no deadlock, because
  \field{LoudSkip} sets \field{diagnosticEmitted}$'$ =
  \field{ztrue} unconditionally, so the fourth conjunct of the bad
  state is never satisfied on this path. \\
\field{GateAppliedDeterminesAllowlist} &
  \field{GateState}'s own invariant, \field{allowlistActive =
  gateApplied}, holds in every state either variant's operations can
  reach, by construction: it is not asserted by any operation, it is
  part of the type every operation's after-state must satisfy. \\
\field{WritePassthroughWhenUngated} &
  Corollary of \field{GateAppliedDeterminesAllowlist} and
  \field{WriteAllowedNoAllowlist}'s missing \field{targetInAllowlist?}
  guard: whenever \field{gateApplied = zfalse} --- for any reason,
  including \field{MuteSkip} firing on a spawn where
  \field{verifierOfOpenMission = ztrue} --- \field{allowlistActive =
  zfalse} and \field{WriteAllowedNoAllowlist} is enabled for every
  \field{Write}/\field{Edit} tool call this spawn makes, with no
  further condition. This is the composition fact
  \S\ref{sec:state}'s \field{allowlistActive} paragraph states: the
  frozen-evaluator gate and write-set enforcement are disabled
  together, by one shared cause. \\
\field{CauseIndistinguishable} &
  \field{MuteSkip}'s guard and its four assignment predicates
  mention \field{channel?} and \field{decided}; none mentions
  \field{cause?}. For any before-state and any \field{c1, c2 :
  AbsenceCause}, firing \field{MuteSkip} with \field{cause? = c1}
  and with \field{cause? = c2} from the same before-state produces
  identical after-states. None of \field{causeTierAAdHoc},
  \field{causePreToolUseMisdispatch}, or
  \field{causePreToolUseNotInstalled} (\S\ref{sec:causes}) is
  distinguishable at this gate --- the channel is a single string;
  presence or absence is the only bit it carries. This holds equally
  for the corrected \field{LoudSkip} in \S\ref{sec:validation}: its
  diagnostic can say a mission binding was absent, never which of the
  three absences produced it. \\
\bottomrule
\end{tabularx}
\end{center}

Two further facts, stated as prose rather than as table rows because
neither is a claim ProB checks or a claim a single predicate states:

\begin{itemize}
  \item \field{GroundTruthFrozen}: every operation in
    \S\ref{sec:decision}--\ref{sec:writeops} carries
    \field{verifierOfOpenMission}$'$ = \field{verifierOfOpenMission}
    (the three decision operations explicitly; the three write
    operations via \field{\Xi GateState}). Ground truth, once fixed
    at \field{Init}, is invariant across every operation this
    document models --- matching \fpath{spec-mission-lifecycle.tex}'s
    \field{EvaluatorHashFrozen} in form and in the reason for stating
    it: the trust anchor a gate reasons about must not itself move
    while the gate is deciding.
  \item \field{SubagentStartOneShot}: \field{MuteSkip},
    \field{NotAVerifier}, and \field{VerifierGateHit} all guard on
    \field{decided = zfalse} and all set \field{decided}$'$ =
    \field{ztrue}. Once one of the three has fired, none of the three
    --- nor any future replacement operation occupying this same
    role --- is enabled again for this spawn, matching
    \field{HandleSubagentStartWithDeps} running exactly once per
    spawn.
\end{itemize}

% ===================================================================
\section{Validation: is silent-disabled reachable?}
\label{sec:validation}
% ===================================================================

The mission brief for this document asks for two things: a recorded
counter-example showing today's \field{MuteSkip} semantics reach the
silent-disabled state, and a demonstration that a corrected,
diagnostic-only variant does not. Following
\fpath{spec-mission-lifecycle.tex}'s own validation technique
exactly, both demonstrations use a state schema stripped to only
what the property needs, checked by reachability/deadlock rather than
by an invariant baked into the schema the buggy operation is typed
against.

\subsection{Errata-in-advance: why the property is not a
\field{GateState} invariant}
\label{sec:errata}

An invariant reading, informally, ``\field{verifierOfOpenMission =
ztrue} $\land$ \field{decided = ztrue} $\implies$ \field{gateApplied
= ztrue} $\lor$ \field{diagnosticEmitted = ztrue}'' is exactly the
safety property this document exists to check. It is deliberately
\emph{absent} from \field{GateState} (\S\ref{sec:state}), for the
reason \fpath{spec-mission-lifecycle.tex} \S6.2 gives in full and
this document will not re-derive at the same length: $\Delta S$
abbreviates $S \land S'$, so had that predicate been part of
\field{GateState}'s own \where-clause, \field{MuteSkip}'s after-state
would have to satisfy it too --- and \field{MuteSkip}'s own frame
conjuncts (\field{gateApplied}$'$ = \field{zfalse},
\field{diagnosticEmitted}$'$ = \field{zfalse}) make that
satisfiable only when \field{verifierOfOpenMission = zfalse}. Firing
\field{MuteSkip} with \field{verifierOfOpenMission = ztrue} would
then have \emph{no valid after-state at all} --- not a state ProB
flags as unsafe, a transition ProB reports as impossible, because the
schema calculus made it impossible by construction before any model
checker ran. A ``0 violations'' report from that version would be a
false negative wearing the shape of a clean bill of health, exactly
as \fpath{spec-mission-lifecycle.tex}'s own \field{AbandonFailOpen}
errata found. The property is checked instead as a reachability
question against a schema, \field{GateU} below, that carries no such
invariant, using \field{IdleU} as the self-loop that is disabled at,
and only at, the one state under test.

\subsection{The decoupled model}

Built and run as two separate, self-contained documents ---
\fpath{.tmp/spec-hook-gates-mute.tex} and
\fpath{.tmp/spec-hook-gates-loud.tex}, scratch files under this
repo's \fpath{.tmp/} convention, neither committed nor part of this
file's own checked namespace, for the same reason
\fpath{spec-mission-lifecycle.tex} keeps its two decoupled documents
out of its own namespace: to avoid \S\ref{sec:errata}'s trap and to
avoid redeclaring \field{MuteSkip} a second time in one namespace.
Both share:

\begin{verbatim}
GateU
  verifierOfOpenMission : ZBOOL
  decided : ZBOOL
  gateApplied : ZBOOL
  diagnosticEmitted : ZBOOL

Init
  GateU'
where
  decided' = zfalse
  gateApplied' = zfalse
  diagnosticEmitted' = zfalse

IdleU
  Xi GateU
where
  not (decided = ztrue and verifierOfOpenMission = ztrue
       and gateApplied = zfalse and diagnosticEmitted = zfalse)
\end{verbatim}

\field{IdleU} is a self-loop enabled everywhere except the one state
this demonstration exists to find: decided, a real verifier of a real
open mission, and neither gated nor diagnosed. That converts the
safety property into a question ProB already answers without a
custom goal: if the operation under test can reach that state,
\emph{nothing} is enabled there once it stops being reachable via a
further decision (the decision operations guard on \field{decided =
zfalse} and have already fired) --- deadlock --- and if the corrected
operation is used instead, the state is never reached, and the full
space is deadlock-free.

The first document adds \field{MuteSkip} (\S\ref{sec:decision},
retyped against \field{GateU}, with \field{channel?} and
\field{cause?} dropped since \field{GateU} carries no channel to
range over --- the operation's predicate never used them either):

\begin{verbatim}
MuteSkip
  Delta GateU
where
  decided = zfalse
  decided' = ztrue
  gateApplied' = zfalse
  diagnosticEmitted' = zfalse
  verifierOfOpenMission' = verifierOfOpenMission
\end{verbatim}

\texttt{/z-spec:check} passes with 0 errors. \texttt{/z-spec:model\_check}
explores 3 states and 7 transitions and reports a deadlock, with this
counter-example trace:

\begin{verbatim}
1  INITIALISATION  (verifierOfOpenMission := ztrue, decided := zfalse,
                     gateApplied := zfalse, diagnosticEmitted := zfalse)
2  MuteSkip         decided: zfalse -> ztrue
                     (gateApplied' = zfalse, diagnosticEmitted' = zfalse,
                      verifierOfOpenMission unchanged = ztrue)
   -- deadlock: no operation enabled --
   decided = ztrue, verifierOfOpenMission = ztrue,
   gateApplied = zfalse, diagnosticEmitted = zfalse
\end{verbatim}

This is the fail-silent bug, reached in two steps: a spawn is, in
fact, the evaluator of an open mission, and the SubagentStart hook
receives an empty \field{MISSION\_ID} for that spawn --- for any of
the three causes in \S\ref{sec:causes}, indistinguishably --- so the
gate that exists to bind the frozen-evaluator context to exactly this
spawn does nothing, and says nothing.

The second document is identical except \field{LoudSkip}
(\S\ref{sec:decision}, retyped against \field{GateU}) replaces
\field{MuteSkip}:

\begin{verbatim}
LoudSkip
  Delta GateU
where
  decided = zfalse
  decided' = ztrue
  gateApplied' = zfalse
  diagnosticEmitted' = ztrue
  verifierOfOpenMission' = verifierOfOpenMission
\end{verbatim}

\texttt{/z-spec:check} passes with 0 errors.
\texttt{/z-spec:model\_check} explores 4 states and 8 transitions and
reports \textbf{no} deadlock, \textbf{no} counter-example, and ALL
states visited --- the state \field{IdleU} excludes is simply never
reached, because \field{LoudSkip} sets
\field{diagnosticEmitted}$'$ = \field{ztrue} unconditionally rather
than by a guard that could fail to fire.

Note on \texttt{-cbc\_deadlock}: probcli's constraint-based deadlock
search, run separately from \texttt{-model\_check}, reports a
deadlocking valuation for \emph{both} documents --- the same
\field{decided = ztrue, verifierOfOpenMission = ztrue, gateApplied =
zfalse, diagnosticEmitted = zfalse} state, found by searching
\field{GateU}'s type directly rather than the states \field{Init} and
the modelled operations actually reach. That state is real in the
type (\field{GateU} carries no invariant forbidding it, deliberately,
per \S\ref{sec:errata}) and unreachable in the \field{LoudSkip}
document (no operation's after-state ever produces it). Constraint-
based deadlock checking does not respect reachability from
\field{Init}; exhaustive \texttt{-model\_check}, which does, is the
correct oracle for a reachability property and is the one this
section's verdicts are drawn from.

\subsection{Tool run}
\label{sec:probrun}

\begin{center}
\begin{tabularx}{\textwidth}{@{}lccl@{}}
\toprule
\textbf{Model} & \textbf{States} & \textbf{Transitions} &
  \textbf{Result} \\
\midrule
Main document (\S\ref{sec:state}--\ref{sec:ops}), 3 decision ops
  $+$ 3 write ops & 6 & 24 & 0 violations, no deadlock \\
Decoupled, \field{MuteSkip} & 3 & 7 & deadlock at
  $\field{decided} = \field{ztrue} \land
  \field{verifierOfOpenMission} = \field{ztrue} \land
  \field{gateApplied} = \field{zfalse} \land
  \field{diagnosticEmitted} = \field{zfalse}$ \\
Decoupled, \field{LoudSkip} & 4 & 8 & clean, all states visited \\
\bottomrule
\end{tabularx}
\end{center}

fuzz's \texttt{/z-spec:check} type-checks this document with 0 errors.
All \texttt{model\_check}/\texttt{animate} runs reported above,
including the two decoupled documents, were run against \texttt{.tex}
files --- an earlier revision of the decoupled documents used an
\field{Init} schema taking \field{verifierOfOpenMission?} as an
explicit input, which probcli's Z-to-B translator rejects for an
\field{INITIALISATION} operation (\field{identifier\_not\_found}: an
initialisation schema has nothing to supply an input at machine
start). \field{verifierOfOpenMission}$'$ is left unconstrained
instead, exactly as \S\ref{sec:init} does in the main document; ProB
enumerates both truth values as separate initial states, which is
what the counter-example trace above reports as
\field{verifierOfOpenMission := ztrue} at step 1 rather than as an
input to step 1.

% ===================================================================
\section{Traceability}
\label{sec:trace}
% ===================================================================

\begin{center}
\begin{tabularx}{\textwidth}{@{}llX@{}}
\toprule
\textbf{Z construct} & \textbf{Go source} & \textbf{Note} \\
\midrule
\field{MissionIDChannel} & \field{os.Getenv("MISSION\_ID")}
  (subagent\_start.go:128) & explicit \field{midAbsent} where Go has
  only \field{""} \\
\field{AbsenceCause} & pretooluse\_dispatch.go:80--132, 894 &
  the three real events that all reach the gate as
  \field{channel? = midAbsent} \\
\field{GateState} & \field{SubagentStartDeps}/
  \field{checkVerifierHash} (subagent\_start.go:51--72, 585--652) $+$
  the \field{ETHOS\_VERIFIER\_ALLOWLIST} env var
  (pretooluse.go:110) & spans both files; see
  \S\ref{sec:onespawn} \\
\field{Init} & the instant before
  \field{HandleSubagentStartWithDeps} makes its one decision & \\
\field{MuteSkip} & \field{checkVerifierHash}
  (subagent\_start.go:599--605) & \emph{current} code; see
  \S\ref{sec:validation} \\
\field{NotAVerifier} & \field{checkVerifierHash}
  (subagent\_start.go:611--621) & the ethos-z69l fix \\
\field{VerifierGateHit} & \field{checkVerifierHash}
  (subagent\_start.go:611--651) & excludes the hash-mismatch refusal,
  \S\ref{sec:nongoals} \\
\field{WriteAllowedNoAllowlist} & \field{HandlePreToolUse}
  (pretooluse.go:110--113) & \\
\field{WriteAllowedInAllowlist}, \field{WriteDenied} &
  \field{HandlePreToolUse}/\field{pathAllowed}
  (pretooluse.go:119--128) & \field{targetInAllowlist?} abstracts
  \field{splitAllowlist}/\field{pathAllowed}, \S\ref{sec:nongoals} \\
\field{LoudSkip} & not present in current code; proposed corrected
  variant, matching the sibling empty-hash warn already at
  subagent\_start.go:624--635 & modelled only for
  \S\ref{sec:validation}; bead ethos-mr1z \\
\bottomrule
\end{tabularx}
\end{center}

\end{document}
